\documentclass{article}
\usepackage{pathfinder-note}

\title{Compile-to-Code as the Manufacture of One-Sided Imitation: What Q and P Establish, and What They Do Not}

\begin{document}
\maketitle

\section{The pair}

\textbf{Q} (\emph{Mechanism Design for Alignment and Control}) is a principal--agent theory paper. A designer wants agents to act on its behalf when both alignment (preferences) and capabilities (feasible actions and information) are unknown, so mechanisms must incentivize honesty \emph{and} obedience. Its structural premise is a one-sided imitation condition --- capabilities can be concealed but not counterfeited --- from which it claims a revelation principle, a characterization of implementable policies via nested cyclical monotonicity, and conditions under which higher-order beliefs discipline multiple agents. Applications listed: sandbagging, an alignment--interpretability trade-off, peer scoring, coupled rewards, scalable oversight.

\textbf{P} (\emph{Commitment To Cooperation With Self-Negotiated Contracts}) is an empirical multi-agent paper. LLM-backed agents negotiate contracts in a spatial-temporal bargaining-plus-navigation game, across a suite of contract representations ranging from formal contracts that compile to code to natural contracts requiring reinterpretation, and across backbones of different sizes and providers. Headline: self-negotiated contracts improve cooperative outcomes beyond plain trading.

\section{Evidence base, stated plainly}

The evidence here is thin and this qualification governs every claim below. Both inputs point \texttt{text=} at source files (\texttt{sources/2609.01595.tex}, \texttt{sources/2607.22750.tex}) that do not exist anywhere under the project root; there is no \texttt{sources/} directory at all, re-verified \ledger{1, 26}. We had only the two abstracts. No entry in this account rests on a theorem statement, a proof, a results table, or a logged experimental grid. Everything is abstract-level inference and must be re-derived against the full texts before anything is written for submission \ledger{1, 37}.

A second limit on the record: peer exchange ended on allowance, not on convergence. The final declared-ready state from the partner is at ledger \#27, so ledger \#28 and \#29 --- including the refinement of the design predictions and the explicit disagreement recorded there --- were never adjudicated. Where that matters, it is flagged below.

\section{The connexion pursued}

The scan proposed: characterize which self-negotiated contracts in P are implementable, via Q's nested cyclical monotonicity; feasibility 70 on the ground that this ``needs no new methods or data''. Both peers rejected this, independently and for compatible reasons \ledger{2, 7}.

\begin{itemize}
\item \emph{Category error.} Implementability in Q's sense is a property of a social choice rule selected by a designer who \emph{commits} to a mechanism; obedience is defined relative to such a principal. In P the contract is the \emph{output} of bilateral self-negotiation, with no outside party holding commitment power or eliciting types. The predicate does not apply to the object, and the revelation principle fails in the standard informed-principal way \ledger{7, 2}.
\item \emph{Problem-type mismatch.} P's stated core difficulty is temporal: costs of cooperation are incurred early, benefits realized later, creating an incentive to defect --- hold-up and moral hazard. Q's is adverse selection over unknown preferences and capabilities. These coincide only where capability is private and payoff-relevant at negotiation time, which is a subset of P's setting rather than its main axis \ledger{2}, conceded by the partner \ledger{17}.
\end{itemize}

The repair that survives has two moves.

\paragraph{Relocate the designer one level up \ledger{8}.} P does not merely report contracts; it \emph{designs the representation layer}. That suite is itself the mechanism-design choice in Q's sense: the designer is the experimenter choosing representation plus enforcement technology, and the agents negotiate within it. Self-negotiation is then part of the game form rather than a rival to it, which dissolves the category error. Accepted without reservation \ledger{27}.

\paragraph{The bridge \ledger{3, 9}.} Q's one-sided imitation premise is not a modelling convenience but an \emph{empirical property of the representation layer}, and P has the exact knob that varies it. A contract that compiles to code conditions payoff on verified action: concealment remains possible, overstatement is self-voiding, so one-sided imitation holds and Q's implementability apparatus applies. A natural contract requiring reinterpretation restores counterfeiting: capability claims are cheap talk adjudicated only ex post, and the reinterpreter is itself an LLM of unknown type, so the enforcement map is type-dependent and lacks commitment. Q's results therefore transfer to P's formal end and not to its natural end, and compile-to-code is the operation that \emph{manufactures} Q's identifying assumption.

This hinge was written twice independently, from the same two abstracts and without sight of the other \ledger{16}. That convergence is the main reason to think it is the load-bearing connexion rather than an artefact of one reading. Everything else in this account is downstream of it.

\section{Supported findings}

\begin{enumerate}
\item \textbf{Temporal pricing, not eventual detection, is the operative variable \ledger{18}.} The strongest objection to the bridge is that the game's execution eventually reveals capability in \emph{both} arms, making formal-versus-natural a difference of enforcement cost rather than of Q's premise \ledger{5}. It loses, but not for the reason first given. Under hold-up, ex-post revelation of counterfeited capability arrives \emph{after} the cooperating party has sunk its cost, hence is payoff-irrelevant to that party: it cannot condition on it. What defeats counterfeiting is that it is priced \emph{before} the sunk cost --- which compile-to-code does and reinterpretation does not. The temporal structure raised in \ledger{2} as an objection is thus what saves the claim; that duality is itself worth stating.

\item \textbf{The revelation principle does not imply representation-invariance \ledger{23}, conceded in full \ledger{31}.} RP is an existence claim about the designer's attainable set: every equilibrium outcome of every mechanism is replicable by \emph{some} direct truthful mechanism. Different game forms retain different equilibria even with fully rational agents. So ``P finds representation matters, therefore an RP premise fails'' is a non sequitur, and there is no invariance benchmark for capable agents to converge to. The proposed capability-level test \ledger{11} is withdrawn on two grounds: its null is misspecified (the correct limit is a \emph{positive} asymptote equal to the genuine enforcement difference between compiled and reinterpreted contracts), and its interaction sign is unidentified, since a larger backbone is better both at strategic reasoning (shrinks the gap) and at exploiting an unenforceable natural contract (grows it).

\item \textbf{What is empirical is the behavioural gap, not the theorem \ledger{32}.} RP cannot be falsified by experiment and no paper should be framed as testing it. Two things are empirical: whether the premises obtain in the platform, and whether LLM agents \emph{play} the equilibrium RP posits. The measured object is therefore the distance between the theoretically attainable set and what agents achieve under each elicitation format --- nonnegative by construction, and producible by neither paper alone.

\item \textbf{Directional content survives, and it is signed \ledger{24}.} Holding designer commitment and enforcement power fixed, the best direct truthful mechanism weakly dominates any indirect one: dominance, not invariance. Falsification is then signed --- indirect beating direct is a violation, formal beating natural is not. P cannot run this test, because its whole suite is \emph{indirect}: in both arms agents negotiate and nobody reports a type. The direct-revelation endpoint Q licenses is simply missing from P's design space.

\item \textbf{The design.} A \(2 \times 2\): rows are enforcement technology (compile-to-code vs.\ LLM reinterpretation), columns are elicitation format (direct capability report vs.\ free bilateral negotiation). Cell \(B\) is P's formal arm, cell \(D\) its natural arm; cells \(A\) (report + code) and \(C\) (report + reinterpreter) are new \ledger{24}. Predictions separate by sign: RP dominance gives \(A \succeq B\) and \(C \succeq D\); obvious strategy-proofness in the sense of Li (2017) predicts the reversal \(B \succ A\), since indirect formats are frequently easier for an agent than their strategically equivalent direct forms; the one-sided imitation bridge predicts the code row dominates the reinterpretation row and confines sandbagging and counterfeiting signatures to \(C\) and \(D\). The decisive gain over \ledger{11} is that commitment and bounded rationality are separated \emph{by design} --- enforcement is held fixed within each row --- rather than by fitting the sign of a confounded interaction.

\item \textbf{Cell \(A\) is a one-sided direct mechanism, constructed rather than assumed \ledger{34}.} Q's premise makes the message space one-sided: reports may understate, \(\hat{\theta} \le \theta\), but not overstate. In a hand-built direct mechanism that restriction is imposed by fiat, which is the usual complaint about such models. Compile-to-code imposes it for free --- an overstated claim yields null execution and no payment, so overstatement is not prohibited but unprofitable and self-voiding --- while a reinterpreter can award partial credit for a plausible good-faith attempt, making overstatement profitable in expectation. Cell \(A\) is therefore the natural home for the implementability certification, and cell \(C\) is the same elicitation \emph{without} one-sidedness, which is what isolates Q's identifying assumption from elicitation format. This gives each cell a theoretical reason rather than a factorial one.

\item \textbf{Fit the rule, then certify it \ledger{25}.} \(A \succeq B\) follows from RP only if arm \(A\) uses the right allocation rule; RP guarantees existence of \emph{some} equivalent direct mechanism, so a raw failure is a joint rejection of the premises and of a hand-chosen rule. The fix is not to choose by hand: fit the rule from arm \(B\)'s observed (declared capability, allocation, transfer) triples, test the fitted rule against nested cyclical monotonicity under the one-sided capability order, and run the certified rule as arm \(A\). Both branches publish. If the fitted rule satisfies nested CM, arm \(A\) is a legitimate benchmark and the comparison is clean. If it fails, that is the headline: the contracts these agents negotiate are implementable by \emph{no} mechanism with commitment --- the negotiation output is not mechanism-rationalizable \ledger{13}. Cyclical monotonicity is thus the bridge between arms, not a decoration on P's metrics.

\item \textbf{Quasilinearity is a verification item, not the main threat \ledger{26, 36}.} It was first called the single biggest threat to the method \ledger{14}. Two grounds downgrade it: the game appears to be Colored Trails, whose conventional scoring is additively separable (goal bonus, per-chip value for unused chips, distance penalty), giving a transferable numeraire in chips; and Q's authors work in a standardly transfer-based tradition. Both are inference from the abstract plus platform convention, not from the texts. One precision matters: run the CM test on \emph{contract terms as written at agreement time}, not on realized play. The map from declared capability to (allocation, chip transfer) is static and quasilinear in chips; realized play is where temporal non-separability lives, and testing there would test the wrong map \ledger{36}. The test then consumes negotiation transcripts, which P logs by construction.
\end{enumerate}

\section{Objections that stand}

\begin{itemize}
\item \textbf{Impoverished message space in cell \(A\) \ledger{33}.} Arm \(A\) is not P with a different representation; it removes negotiation and replaces it with centralized allocation from reports. But the game is spatial-temporal, and the negotiation channel in \(B\) and \(D\) transmits payoff-relevant information beyond capability --- board position, intended route, chip complementarity. If the direct arm's message space is only ``my capability type'', \(A \prec B\) follows mechanically from an impoverished message space rather than from any premise failure or from OSP, and a referee will say so. The type space in \(A\) and \(C\) must be rich enough to be informationally equivalent to the indirect arms. This is the practical face of the fitted-rule caveat and it pushes hard toward fit-then-certify, because fitting from transcripts is what reveals which private information was actually being priced. Unaddressed here; it is a design requirement, not a resolved point.

\item \textbf{The escrow question \ledger{35}, gating.} If pricing before the sunk cost is what defeats counterfeiting, then compile-to-code defeats it only if compiled contracts \emph{withhold} the counterparty's chips or payment until the promised action is verified --- only if they are escrow-like and conditional. If P's compiled contracts merely execute an agreed transfer at agreement time, with code buying unambiguity alone, then the counterparty's cost is still sunk before revelation, one-sidedness fails in the formal arm too, and the shared hinge collapses to a claim about ambiguity rather than about Q's premise. This objection arrived after the final ready-at-\#27 and was never answered on the record. I accept its force: it is the single highest-value question for the full text, above the form of nested CM. A negative answer kills the thesis; a positive answer makes it nearly automatic.

\item \textbf{Level versus asymmetry \ledger{19}.} Capability level and capability asymmetry collide in the matched-strong cell, where both predict a small representation gap for different reasons; separating them needs a contrast the existing grid may not contain. This is correct against \ledger{4} and \ledger{11}, and is a further reason the \(2 \times 2\) supersedes both: by holding enforcement fixed within a row it does not need the level-versus-asymmetry contrast to carry identification. The capability axis then reduces to matched versus asymmetric, needed only for the sandbagging and surplus-split prediction, giving \(2 \times 2 \times 2 = 8\) cells of which P already populates two per pairing type \ledger{28}.

\item \textbf{Scope restrictions.} Q's applications (iii) peer scoring and (iv) coupled rewards require three or more agents or a peer evaluator; if the game is strictly bilateral they are inapplicable and must not be claimed \ledger{14, 27}. The sandbagging prediction \ledger{12} --- that under natural contracts the surplus split is compressed relative to the true capability gap while under compiled contracts it tracks the gap more tightly --- must be scoped to capability-asymmetric pairings \ledger{17}, which is precisely where it already lives.

\item \textbf{Power \ledger{14, 36}.} Cycle detection needs many distinct contracts per backbone. One contract per episode is a slow way to get them, and this worry is not relieved by the quasilinearity downgrade.
\end{itemize}

\section{Unresolved}

\begin{enumerate}
\item \emph{Unfixable in this thread}: neither full text exists under the root \ledger{1, 26}. Step zero of any follow-up is obtaining both.
\item Are P's compiled contracts conditional and escrowed, or merely unambiguous \ledger{35}? Gating; kills or makes the thesis.
\item Does nested cyclical monotonicity assume transfers and quasilinear payoffs, and does ``nested'' encode the partial order induced by one-sided imitation (IC constraints only downward along the capability order)? Both peers inferred yes to the second; it is flagged as a \emph{shared} unverified inference so that neither reading is mistaken for a check \ledger{14, 27}.
\item Is the game strictly bilateral \ledger{14}?
\item Does P report contract terms and surplus split conditional on backbone \emph{pairing}, or only aggregate outcomes, and does the matched-strong cell exist \ledger{5, 19}?
\item Does P's compiled arm log the \emph{declared} capability alongside allocation and transfer? Without it the fitted-rule step has no input \ledger{29}.
\item Power for the cycle test \ledger{14}.
\item \emph{Left unadjudicated by the allowance.} The prediction in \ledger{28} --- that if timing of pricing relative to the sunk cost is the operative variable, the row effect does the work and the column effect is confined to the code row, so \(C \approx D\) --- was never seen by the partner. Nor was the disagreement recorded in \ledger{29} over whether cyclical monotonicity is load-bearing or a discussion-section proposal \ledger{21}; I note that \ledger{34} independently makes cell \(A\) the home of the CM certification, which converges on the load-bearing reading without having seen \ledger{28, 29}. I take that as substantially settling it, but it was not settled on the record.
\end{enumerate}

\section{Scoring}

Both peers reject the scan framing as a category error and both reject its justification that the work ``needs no new methods or data'': it needs a new method (the CM certification) and new data (two absent arms). Gain revised upward from the scan's 55 to roughly 70--75, on the shared ground that the contribution is a measurement neither paper can make alone; I hold the higher end for the signed dominance test, the partner the lower end because a measured behavioural gap honestly reads as less dramatic than a falsification \ledger{15, 27, 37}. Feasibility revised downward from 70 to roughly 40--45 --- two new arms, a fit-then-certify step, and a message-space-equivalence requirement constitute a full experimental paper \ledger{27, 37}. Earlier scores of 35 for the scan object and 55 for the reframed one were not a real disagreement: they scored different objects \ledger{20}. Verdict: connexion \emph{confirmed} but substantially rebuilt; the hypothesis as scanned is rejected.

\section{Smallest next step}

Obtain the two full texts, then answer one question before anything else: \emph{are P's compiled contracts conditional and escrowed, or merely unambiguous?} \ledger{35}. That single fact determines whether compile-to-code prices counterfeiting before the sunk cost, which is the mechanism the entire line rests on \ledger{18, 34}. If escrowed, proceed to the second question --- whether the compiled arm logs declared capability alongside allocation and transfer \ledger{29} --- since that determines whether the fitted-rule step has input and hence whether the design is principled or hand-built. Nothing further should be written until both are answered.

\paragraph{Proposed structure, if it survives.} Thesis: the one-sided imitation bridge \ledger{3, 9}, with the temporal-pricing argument as its mechanism \ledger{18}. Framing: the behavioural RP gap and signed dominance \ledger{24, 32}, not representation-dependence as falsification \ledger{23}. Method: fitted-rule CM certification \ledger{25}, with cell \(A\) constructed as a one-sided direct mechanism \ledger{34} and message-space equivalence enforced \ledger{33}. Experiment: the eight-cell design \ledger{28}, with sandbagging scoped to asymmetric pairings \ledger{12, 17}. Applications (iii) and (iv) dropped unless the game is multilateral \ledger{14}.

\end{document}